\documentclass[12pt]{article}
\usepackage[margin=1in]{geometry}
\usepackage{amsmath}
\usepackage{amssymb}
\usepackage{amsthm}
\usepackage{enumitem}
\usepackage{xcolor}
\usepackage{pagecolor}
\usepackage{marginnote}
\usepackage{fancyhdr}
\usepackage{bm}
\usepackage{etoolbox}
\usepackage{mdframed}

\definecolor{cream}{HTML}{faf7f1}
\pagecolor{cream}

\newcommand{\defn}[1]{\textbf{#1}}
\newcommand{\defeq}{\mathrel{\vcenter{\baselineskip0.5ex \lineskiplimit0pt
                     \hbox{\scriptsize.}\hbox{\scriptsize.}}}=}
\newcommand{\T}{\top}
\newcommand{\F}{\bot}
\newcommand{\Pow}{\mathbb{P}}
\newcommand{\suc}{\mathfrak{suc}}
\newcommand{\N}{\mathbb{N}}

\newtheoremstyle{proofstyle}
  {}{}{}{}{\bfseries}{.}{ }{}
\theoremstyle{proofstyle}
\newtheorem*{proof*}{Proof}

\newcommand{\sidenote}[1]{\marginnote{\small\itshape #1}}

\makeatletter
\renewcommand{\maketitle}{%
  \begin{flushleft}
    {\Large\bfseries CS173: Discrete Structures}\\[0.3em]
    {\Large\bfseries Problem Set 5}\\[0.5em]
    {\normalsize Mohammed Al Salhi}
  \end{flushleft}
  \vspace{1em}
}
\makeatother

\AtBeginDocument{\mathversion{bold}}

\begin{document}

\maketitle

\begin{enumerate}[left=0pt, itemsep=2em, parsep=1em]

%% ---------------------------------------------------------------
%% Problem 1
%% ---------------------------------------------------------------
\item We say that a set $x$ is \defn{$\in$-transitive} $:\Leftrightarrow \forall y(y \in x \Rightarrow \forall z(z \in y \Rightarrow z \in x))$.

Prove that every natural number is $\in$-transitive.

\begin{proof*}
We will show $(\forall n \in \N)(n \text{ is } {\in}\text{-transitive})$ by weak induction.

\medskip
\textbf{Basis Step:}
We must show $0$ is $\in$-transitive, i.e., $\forall y(y \in 0 \Rightarrow \forall z(z \in y \Rightarrow z \in 0))$. Since $0 = \emptyset$, there is no $y$ such that $y \in 0$. Therefore any statement $\forall y(y \in 0 \Rightarrow p)$ where p is an arbitrary proposition holds.

\medskip
\textbf{Inductive Step:}
Let $k \in \N$, and assume $k$ is $\in$-transitive. We will show $\suc(k)$ is $\in$-transitive.

Recall $\suc(k) = k \cup \{k\}$ by definition. Let $y$ be arbitrary and suppose $y \in \suc(k)$. Then $y \in k \cup \{k\}$, so either $y \in k$ or $y \in \{k\}$, meaning either $y \in k$ or $y = k$.

Now let $z$ be arbitrary and suppose $z \in y$. We consider two cases.

\medskip
\textit{Case 1:} Suppose $y \in k$. Since $k$ is $\in$-transitive by the inductive hypothesis and $y \in k$ and $z \in y$, we conclude $z \in k$. Since $k \subseteq k \cup \{k\} = \suc(k)$, we have $z \in \suc(k)$.

\medskip
\textit{Case 2:} Suppose $y = k$. Then $z \in y = k$, so $z \in k \subseteq k \cup \{k\} = \suc(k)$, giving us $z \in \suc(k)$.

\medskip
In either case, $z \in \suc(k)$. Since $y$ and $z$ were arbitrary, $\suc(k)$ is $\in$-transitive.

Therefore, we conclude $(\forall n \in \N)(n \text{ is } {\in}\text{-transitive})$.
\end{proof*}

\newpage

%% ---------------------------------------------------------------
%% Problem 2
%% ---------------------------------------------------------------
\item We will work up to a proof of commutativity of addition.

\begin{enumerate}[label=(\alph*), itemsep=1.5em, leftmargin=2em]

  \item Prove $(\forall x, y \in \N)(\suc(x + y) = \suc(x) + y)$.

  \begin{proof*}
  Let $x \in \N$. We will prove $(\forall y \in \N)(\suc(x + y) = \suc(x) + y)$ by weak induction.

  \medskip
  \textbf{Basis Step:}
  We observe the following.
  \[
    \suc(x) + 0 = \suc(x) \tag{by definition of $+$}
  \]
  \[
    \suc(x + 0) = \suc(x) \tag{by definition of $+$}
  \]
  Therefore $\suc(x + 0) = \suc(x) + 0$.

  \medskip
  \textbf{Inductive Step:}
  Let $k \in \N$, and assume $\suc(x + k) = \suc(x) + k$. We observe the following.
  \begin{align*}
    \suc(x) + \suc(k) &= \suc(\suc(x) + k) &&\text{by definition of }+ \\
                      &= \suc(\suc(x + k)) &&\text{by the inductive hypothesis}\\
                      &= \suc(x + \suc(k)) &&\text{by definition of }+
  \end{align*}
  Therefore $\suc(x + \suc(k)) = \suc(x) + \suc(k)$.

  We can therefore conclude $(\forall y \in \N)(\suc(x + y) = \suc(x) + y)$, and since $x$ was arbitrary, $(\forall x, y \in \N)(\suc(x + y) = \suc(x) + y)$.
  \end{proof*}

  \item Prove $(\forall x, y \in \N)(x + y = y + x)$.

  \begin{proof*}
  Let $x \in \N$. We will prove $(\forall y \in \N)(x + y = y + x)$ by weak induction.

  \medskip
  \textbf{Basis Step:}
  We observe the following.
  \begin{align*}
    x + 0 &= x &&\text{by Theorem: Zero is the Additive Identity} \\
          &= 0 + x &&\text{by Theorem: Zero is the Additive Identity}
  \end{align*}
  Therefore $x + 0 = 0 + x$ as desired.

  \medskip
  \textbf{Inductive Step:}
  Let $k \in \N$, and assume $x + k = k + x$. We observe the following.
  \begin{align*}
    x + \suc(k) &= \suc(x + k) &&\text{by definition of }+ \\
                &= \suc(k + x) &&\text{by the inductive hypothesis} \\
                &= \suc(k) + x &&\text{by part (a)}
  \end{align*}
  Therefore $x + \suc(k) = \suc(k) + x$ as desired.

  We can therefore conclude $(\forall y \in \N)(x + y = y + x)$, and since $x$ was arbitrary, $(\forall x, y \in \N)(x + y = y + x)$.
  \end{proof*}

\end{enumerate}

\newpage

%% ---------------------------------------------------------------
%% Problem 3
%% ---------------------------------------------------------------
\item Show that $(\forall x, y \in \N)(x < y \Rightarrow (\exists n \in \N)(x + n = y))$.

\begin{proof*}
Let $x \in \N$. We will prove $(\forall y \in \N)(x < y \Rightarrow (\exists n \in \N)(x + n = y))$ by weak induction.

\medskip
\textbf{Basis Step:}
We must show $x < 0 \Rightarrow (\exists n \in \N)(x + n = 0)$. Recall that $x < 0$ means $x \in 0 = \emptyset$ by definition, which is false for any $x$. Therefore the implication $x < 0 \Rightarrow \ldots$ holds vacuously.

\medskip
\textbf{Inductive Step:}
Let $k \in \N$, and assume $x < k \Rightarrow (\exists n \in \N)(x + n = k)$. We must show $x < \suc(k) \Rightarrow (\exists n \in \N)(x + n = \suc(k))$.

Suppose $x < \suc(k)$, meaning $x \in \suc(k) = k \cup \{k\}$. Then either $x \in k$ or $x \in \{k\}$, i.e., either $x < k$ or $x = k$.

\medskip
\textit{Case 1:} Suppose $x < k$. By the inductive hypothesis, there exists $n_0 \in \N$ such that $x + n_0 = k$. Then observe:
\begin{align*}
  x + \suc(n_0) &= \suc(x + n_0) &&\text{by definition of }+ \\
                &= \suc(k)        &&\text{since }x + n_0 = k
\end{align*}
So $\suc(n_0) \in \N$ witnesses $(\exists n \in \N)(x + n = \suc(k))$.

\medskip
\textit{Case 2:} Suppose $x = k$. Then observe:
\begin{align*}
  x + 1 &= \suc(x) &&\text{by Lemma 6.5}\\
        &= \suc(k) &&\text{since }x = k
\end{align*}
So $1 \in \N$ witnesses $(\exists n \in \N)(x + n = \suc(k))$.

\medskip
In either case, $(\exists n \in \N)(x + n = \suc(k))$ holds. We therefore conclude $x < \suc(k) \Rightarrow (\exists n \in \N)(x + n = \suc(k))$.

Since $x$ was arbitrary, $(\forall x, y \in \N)(x < y \Rightarrow (\exists n \in \N)(x + n = y))$ as desired.
\end{proof*}

\newpage

%% ---------------------------------------------------------------
%% Problem 4
%% ---------------------------------------------------------------
\item Prove that $(\forall x, y, z \in \N)(x \cdot (y + z) = (x \cdot y) + (x \cdot z))$.

\begin{proof*}
Let $x, y \in \N$. We will prove $(\forall z \in \N)(x \cdot (y + z) = (x \cdot y) + (x \cdot z))$ by weak induction.

\medskip
\textbf{Basis Step:}
We observe the following.
\begin{align*}
  x \cdot (y + 0) &= x \cdot y &&\text{by definition of }+ \\
                  &= (x \cdot y) + 0 &&\text{by definition of }\cdot \\
                  &= (x \cdot y) + (x \cdot 0) &&\text{by definition of }\cdot
\end{align*}
Therefore $x \cdot (y + 0) = (x \cdot y) + (x \cdot 0)$ as desired.

\medskip
\textbf{Inductive Step:}
Let $k \in \N$, and assume $x \cdot (y + k) = (x \cdot y) + (x \cdot k)$. We observe the following.
\begin{align*}
  x \cdot (y + \suc(k))
    &= x \cdot \suc(y + k) &&\text{by definition of }+ \\
    &= x \cdot (y + k) + x &&\text{by definition of }\cdot \\
    &= ((x \cdot y) + (x \cdot k)) + x &&\text{by the inductive hypothesis} \\
    &= (x \cdot y) + ((x \cdot k) + x) &&\text{by associativity of }+ \\
    &= (x \cdot y) + (x \cdot \suc(k)) &&\text{by definition of }\cdot
\end{align*}
Therefore $x \cdot (y + \suc(k)) = (x \cdot y) + (x \cdot \suc(k))$.

We can therefore conclude $(\forall z \in \N)(x \cdot (y + z) = (x \cdot y) + (x \cdot z))$, and since $x$ and $y$ were arbitrary, $(\forall x, y, z \in \N)(x \cdot (y + z) = (x \cdot y) + (x \cdot z))$.
\end{proof*}

\newpage

%% ---------------------------------------------------------------
%% Problem 5
%% ---------------------------------------------------------------
\item Given a sequence of natural numbers $f : \N \to \N$ and a natural number $a \in \N$, we recursively define the iterated summation of finitely many values of $f$ as follows:
\[
  \sum_{i=a}^{a} f(i) \defeq f(a)
\]
\[
  \sum_{i=a}^{\suc(b)} f(i) \defeq \left(\sum_{i=a}^{b} f(i)\right) + f(\suc(b))
  \quad \text{where } b \in \N \text{ such that } a \leq b
\]
For completeness, we say $\displaystyle\sum_{i=a}^{b} f(i) \defeq 0$ when $b \in \N$ such that $b < a$. For this problem, you may also assume associativity and commutativity of multiplication over $\N$.

Prove that $\displaystyle(\forall n \in \N)\!\left(1 + \sum_{i=0}^{n} 2^i = 2^{n+1}\right)$.

\begin{proof*}
We will prove $(\forall n \in \N)\!\left(1 + \displaystyle\sum_{i=0}^{n} 2^i = 2^{n+1}\right)$ by weak induction.

\medskip
\textbf{Basis Step:}
We observe the following. In the basis step, we need to prove $1 + \displaystyle\sum_{i=0}^{0} 2^i = 2^{0+1}$.
\begin{align*}
  1 + \sum_{i=0}^{0} 2^i
    &= 1 + 2^0 &&\text{by definition of iterated summation} \\
    &= 1 + 1   &&\text{by definition of exponentiation} \\
    &= 2       &&\text{by definition of +} \\
    &= 2^1     &&\text{since }2^1 = 2^{\suc(0)} = 2^0 \cdot 2 = 1 \cdot 2 = 2\\
    &= 2^{0+1} &&\text{since }0 + 1 = 1
\end{align*}
Therefore $1 + \displaystyle\sum_{i=0}^{0} 2^i = 2^{0+1}$ as desired.

\medskip
\textbf{Inductive Step:}
Let $k \in \N$, and assume $1 + \displaystyle\sum_{i=0}^{k} 2^i = 2^{k+1}$. In the inductive step, we must prove $1 + \displaystyle\sum_{i=0}^{\suc(k)} 2^i = 2^{\suc(k)+1}$.
We observe the following.
\begin{align*}
  1 + \sum_{i=0}^{\suc(k)} 2^i
    &= 1 + \left(\sum_{i=0}^{k} 2^i\right) + 2^{\suc(k)}
      &&\text{by definition of iterated summation} \\
    &= \left(1 + \sum_{i=0}^{k} 2^i\right) + 2^{\suc(k)}
      &&\text{by associativity of }+ \\
    &= 2^{k+1} + 2^{\suc(k)}
      &&\text{by the inductive hypothesis} \\
    &= 2^{k+1} + 2^{k+1}
      &&\text{since }(\forall x \in \N)(\suc(x) = x + 1) \\
    &= 2 \cdot 2^{k+1}
      &&\text{by commutativity of }\cdot\text{ and definition of }\cdot \\
    &= 2^{k+1} \cdot 2
      &&\text{by commutativity of }\cdot \\
    &= 2^{\suc(k+1)}
      &&\text{by definition of exponentiation} \\
    &= 2^{(k+1)+1}
      &&\text{since }(\forall x \in \N)(\suc(x) = x + 1) \\
    &= 2^{\suc(k)+1}
      &&\text{by associativity of }+
\end{align*}
Thus, we have shown $1 + \displaystyle\sum_{i=0}^{\suc(k)} 2^i = 2^{\suc(k)+1}$ as required.

Therefore, we can conclude $(\forall n \in \N)\!\left(1 + \displaystyle\sum_{i=0}^{n} 2^i = 2^{n+1}\right)$.
\end{proof*}

\end{enumerate}

\end{document}